% ============================================================
%  A8 - Typografie-Feinschliff: microtype, csquotes, siunitx
%  Kompilieren: pdflatex typografie.tex
% ============================================================
\documentclass[a4paper,11pt]{article}
\usepackage[T1]{fontenc}
\usepackage{lmodern}          % skalierbare Schrift (Voraussetzung für microtype-Expansion)
\usepackage[ngerman]{babel}
\usepackage{booktabs}

% microtype: feinerer Randausgleich und weniger Trennungen.
% Einfach laden, keine weiteren Befehle nötig.
\usepackage{microtype}

% csquotes: \enquote{...} setzt die zur Sprache passenden Anführungszeichen.
\usepackage[autostyle]{csquotes}

% siunitx: Zahlen und Einheiten einheitlich setzen.
%   locale=DE: Dezimalkomma statt Dezimalpunkt in der Ausgabe
\usepackage[locale=DE]{siunitx}

\begin{document}

\section*{Anführungszeichen}
Falsch getippt: "Hallo Welt". Richtig mit csquotes: \enquote{Hallo Welt}.
Auch verschachtelt: \enquote{Er sagte \enquote{Hallo} und ging.}

\section*{Zahlen und Einheiten}
\begin{itemize}
  \item Zahl mit Tausendertrennung: \num{1234567.89}
  \item Zahl mit Einheit: \qty{12.5}{\kilo\metre\per\hour}
  \item Nur die Einheit: \unit{\micro\second}
  \item Bereich: \qtyrange{10}{20}{\celsius}
  \item Speichergröße: \qty{512}{\mega\byte}
\end{itemize}

\section*{Tabelle mit ausgerichteten Dezimalstellen}
Die Spalte \texttt{S} richtet Zahlen am Dezimaltrennzeichen aus.
Text in einer \texttt{S}-Spalte gehört in geschweifte Klammern.

\begin{table}[ht]
  \centering
  \begin{tabular}{l S[table-format=3.2] S[table-format=2.1]}
    \toprule
    Gerät & {Startzeit in \unit{\milli\second}} & {Akku in \unit{\percent}} \\
    \midrule
    Smartphone A & 120.5  & 87.5 \\
    Smartphone B &  98.25 &  9.0 \\
    Tablet C     & 240    & 54.3 \\
    \bottomrule
  \end{tabular}
  \caption{Messwerte, mit \texttt{siunitx} gesetzt}
\end{table}

\section*{Randausgleich}
Dieser Absatz dient als Testtext für \texttt{microtype}. Kommentieren Sie das
Paket in der Präambel aus und vergleichen Sie den rechten Rand und die Anzahl
der Trennungen. Der Unterschied ist klein, summiert sich bei langen Dokumenten
wie einer Bachelorarbeit aber deutlich.

\end{document}
